% Test of side captions (didactic, memoir): the contents of a float with a
% sidecaption go next to the caption (sidecapalign=caption, the default).
% Twoside, captions in the outer margin (didactic's default): the right
% margin on odd pages, the left on even pages.
%
% Build: see the Makefile (pdflatex twice, for the page check).
%
% Expected in the resulting PDF:
%   1  Page 1 (odd): the narrow table TABODD at the right edge of the text
%      block, its caption in the right margin beside it.
%   2  Page 2 (even): the narrow table TABEVEN at the left edge of the text
%      block, its caption in the left margin.
%   3  Page 3 (odd): the narrow image at the right edge, caption right.
%   4  Page 4 (even): an image as wide as the text, which doesn't move; the
%      caption in the left margin.
%   5  Page 5 (odd): a table with \flushscap, same as 1 (TABFLUSH).
%   The captions are at the same height as before, at the top of the
%   contents.
% test-sidecap-off.tex inputs this with sidecapalign=off: then 1, 2, 3
% are at the left of the text block, as memoir places them, while 5, with
% \flushscap, is at the right on its odd page.
\documentclass[a4paper,twoside]{memoir}
\usepackage{graphicx}
\usepackage{lipsum}
\usepackage{didactic}
\begin{document}
\lipsum[1]

\begin{table}[h]
\begin{sidecaption}{A narrow table on an odd page.}[tab:odd]
\begin{tabular}{ll}
TABODD & 1 \\
B & 2 \\
\end{tabular}
\end{sidecaption}
\end{table}

\clearpage
\lipsum[2]

\begin{table}[h]
\begin{sidecaption}{A narrow table on an even page.}[tab:even]
\begin{tabular}{ll}
TABEVEN & 1 \\
B & 2 \\
\end{tabular}
\end{sidecaption}
\end{table}

\clearpage
\lipsum[3]

\begin{figure}[h]
\begin{sidecaption}{A narrow image on an odd page.}[fig:odd]
\includegraphics[width=0.3\linewidth]{example-image-a}
\end{sidecaption}
\end{figure}

\clearpage
\lipsum[4]

\begin{figure}[h]
\begin{sidecaption}{An image as wide as the text, on an even page.}[fig:full]
\includegraphics[width=\linewidth]{example-image-b}
\end{sidecaption}
\end{figure}

\clearpage
\lipsum[5]

\begin{table}[h]
\begin{sidecaption}{A table with flushscap, on an odd page.}[tab:flush]
\flushscap
\begin{tabular}{ll}
TABFLUSH & 1 \\
B & 2 \\
\end{tabular}
\end{sidecaption}
\end{table}
\end{document}
